% !TEX program = pdflatex
% arara: pdflatex: { shell: false }
% arara: pdflatex: { shell: false }
% arara: latexmk: { options: ["-c"] }

\documentclass[a4paper, 10pt]{article}
\let\rnc\renewcommand
\let\nc\newcommand

% ──────────────────────────────────────────────
%  Encodage & langue
% ──────────────────────────────────────────────
\usepackage[T1]{fontenc}
\usepackage[utf8]{inputenc}
\usepackage[french]{babel}

% ──────────────────────────────────────────────
%  Polices
% ──────────────────────────────────────────────
\usepackage{cabin}
\usepackage{inconsolata}
\usepackage{fourier-orns}
\rnc{\familydefault}{\sfdefault}

% ──────────────────────────────────────────────
%  Mise en page
% ──────────────────────────────────────────────
\usepackage[top=1.5cm, bottom=2cm, left=1.5cm, right=1.5cm, headheight=14pt]{geometry}

\usepackage{fancyhdr, lastpage}
\pagestyle{fancy}
\fancyhf{}
\fancyfoot[L]{\small\itshape\CoursClasse~-~\CoursTitre}
\fancyfoot[R]{\thepage~/~\pageref*{LastPage}}
\rnc{\headrulewidth}{0pt}
\rnc{\footrulewidth}{0.4pt}
\fancypagestyle{plain}{%
	\fancyhf{}
	\fancyfoot[L]{\small\itshape\CoursClasse~-~\CoursTitre}
	\fancyfoot[R]{\thepage~/~\pageref*{LastPage}}
	\rnc{\headrulewidth}{0pt}
	\rnc{\footrulewidth}{0.4pt}
}


% ──────────────────────────────────────────────
%  Couleurs
% ──────────────────────────────────────────────
\usepackage[dvipsnames, table]{xcolor}

\definecolor{coulPrimaire}{HTML}{2C3E50}
\definecolor{coulAccent}{HTML}{2980B9}
\definecolor{coulFond}{HTML}{F4F6F8}
\definecolor{coulBordure}{HTML}{AEB6BF}
\definecolor{coulAlerte}{HTML}{C0392B}
\definecolor{coulVert}{HTML}{27AE60}
\definecolor{coulOrange}{HTML}{E67E22}
\definecolor{amethyst}{rgb}{0.6, 0.4, 0.8}
\definecolor{auburn}{rgb}{0.43, 0.21, 0.1}
\definecolor{caribbeangreen}{rgb}{0.0, 0.8, 0.6}
\definecolor{goldenrod}{rgb}{0.85, 0.65, 0.13}
\definecolor{mulberry}{rgb}{0.77, 0.29, 0.55}
\definecolor{tealblue}{rgb}{0.21, 0.46, 0.53}
\definecolor{terracotta}{rgb}{0.89, 0.45, 0.36}
\definecolor{asparagus}{rgb}{0.53, 0.66, 0.42}

% ──────────────────────────────────────────────
%  Colonnes
% ──────────────────────────────────────────────
\usepackage{multicol}
\setlength{\columnsep}{0.5cm}
\setlength{\columnseprule}{0.4pt}
\def\columnseprulecolor{\color{coulBordure}}

% ──────────────────────────────────────────────
%  Maths
% ──────────────────────────────────────────────
\usepackage{amsfonts, amsmath, amssymb, mathtools}
\usepackage{mathrsfs}
\usepackage{yhmath}
\usepackage{xfp}
\usepackage{numprint}
\usepackage{pifont}
\usepackage{relsize}
\usepackage{eurosym}
\mathsurround=2pt

% ──────────────────────────────────────────────
%  TikZ & graphiques
% ──────────────────────────────────────────────
\usepackage{tikz}
\usetikzlibrary{arrows.meta, calc, patterns, decorations.markings,
	shapes.geometric, positioning, backgrounds}
\usepackage{pgfplots}
\pgfplotsset{compat=1.18}
\usepackage{tkz-tab}

% ──────────────────────────────────────────────
%  Figures & tableaux
% ──────────────────────────────────────────────
\usepackage{graphicx}
\usepackage{wrapfig}
\usepackage{float}
\usepackage{booktabs, array, tabularx}
\newcolumntype{R}[1]{>{\raggedleft\arraybackslash}b{#1}}
\newcolumntype{L}[1]{>{\raggedright\arraybackslash}b{#1}}
\newcolumntype{C}[1]{>{\centering\arraybackslash}b{#1}}

% ──────────────────────────────────────────────
%  Code
% ──────────────────────────────────────────────
\usepackage{listings}
\lstset{
	basicstyle=\ttfamily\small,
	breaklines=true,
	frame=none,
	numbers=left,
	numberstyle=\tiny\color{coulBordure},
	commentstyle=\color{coulBordure}\itshape,
	keywordstyle=\color{coulAccent}\bfseries,
	stringstyle=\color{coulAlerte},
	showstringspaces=false,
	xleftmargin=1cm,
}

% ──────────────────────────────────────────────
%  Divers
% ──────────────────────────────────────────────
\usepackage{etoolbox}
\usepackage{fontawesome5}
\usepackage{soul}
\usepackage{contour, ulem}
\rnc{\ULdepth}{3pt}
\rnc{\ULthickness}{1pt}
\contourlength{1pt}
\rnc{\underline}[1]{\uline{\phantom{#1}}\llap{\contour{white}{#1}}}
\usepackage{parskip}
\setlength{\parskip}{0.5mm}
\usepackage{enumitem}
\setlist[itemize]{label=\textbullet, topsep=2pt, itemsep=1pt, parsep=0pt}
\setlist[enumerate]{topsep=2pt, itemsep=1pt, parsep=0pt}

% ──────────────────────────────────────────────
%  tcolorbox
% ──────────────────────────────────────────────
\usepackage{tcolorbox}
\tcbuselibrary{skins, breakable}

\tcbset{
	styleBase/.style={
			enhanced, breakable, arc=2pt, boxrule=0.5pt,
			left=6pt, right=6pt, top=4pt, bottom=4pt,
			fonttitle=\bfseries\small,
		},
	styleLateral/.style={
			enhanced, breakable, frame hidden, opacityback=0,
			left=8pt, right=6pt, top=3pt, bottom=3pt, arc=0pt,
		},
}

% ──────────────────────────────────────────────
%  Liens
% ──────────────────────────────────────────────
\usepackage[
	colorlinks=true,
	linkcolor=coulAccent,
	urlcolor=coulAccent,
	citecolor=coulAccent,
	pdftitle={Fiche d'exercices},
	pdfauthor={},
]{hyperref}

% ──────────────────────────────────────────────
%  Macros générales
% ──────────────────────────────────────────────
\nc{\hligne}{\begin{center}\rule{0.5\textwidth}{1pt}\end{center}}
\nc{\cimg}[2]{\begin{center}\includegraphics[width=#1]{#2}\end{center}}
\nc{\wfig}[2]{\begin{wrapfigure}{r}{#1}\includegraphics[width=#1]{#2}\end{wrapfigure}}
\def\be{\begin{enumerate}}\def\ee{\end{enumerate}}
\def\bi{\begin{itemize}}\def\ei{\end{itemize}}
\let\ms\medskip

% ──────────────────────────────────────────────
%  Macros mathématiques
% ──────────────────────────────────────────────
\nc{\R}{\ensuremath{\mathbb{R}}}
\nc{\D}{\ensuremath{\mathbb{D}}}
\nc{\Q}{\ensuremath{\mathbb{Q}}}
\nc{\Z}{\ensuremath{\mathbb{Z}}}
\nc{\N}{\ensuremath{\mathbb{N}}}
\nc{\C}{\ensuremath{\mathbb{C}}}
\nc{\Cf}{\ensuremath{\mathcal{C}_f}}
\nc{\Cg}{\ensuremath{\mathcal{C}_g}}
\usepackage{esvect}
\rnc{\vec}[1]{\ensuremath{\vv{#1}}}
\nc{\norm}[1]{\ensuremath{\left\lVert#1\right\rVert}}
\nc{\abs}[1]{\ensuremath{\left\lvert#1\right\rvert}}
\nc{\pa}[1]{\ensuremath{\left(#1\right)}}
\nc{\bra}[1]{\ensuremath{\left[#1\right]}}
\nc{\brac}[1]{\ensuremath{\left\{#1\right\}}}
\nc{\pfrac}[2]{\ensuremath{\pa{\frac{#1}{#2}}}}
\nc{\bfrac}[2]{\ensuremath{\bra{\frac{#1}{#2}}}}
\nc{\mat}[1]{\ensuremath{\begin{matrix}#1\end{matrix}}}
\nc{\pmat}[1]{\ensuremath{\pa{\mat{#1}}}}
\nc{\bmat}[1]{\ensuremath{\bra{\mat{#1}}}}
\nc{\smat}[1]{\ensuremath{\begin{smallmatrix}#1\end{smallmatrix}}}
\nc{\psmat}[1]{\ensuremath{\pa{\smat{#1}}}}
\nc{\determinant}[1]{\ensuremath{\left|\mat{#1}\right|}}
\nc{\coord}[2]{\ensuremath{\begin{pmatrix}#1\\#2\end{pmatrix}}}
\nc{\coordl}[2]{\ensuremath{\left(#1~;~#2\right)}}
\nc{\OIJ}{\ensuremath{\left(O;I,J\right)}}
\nc{\vOIJ}{\ensuremath{\left(O;\vec{i},\vec{j}\right)}}
\nc{\eqv}{\Leftrightarrow}
\nc{\rarr}{\rightarrow}\nc{\larr}{\leftarrow}
\nc{\Rarr}{\Rightarrow}\nc{\Larr}{\Leftarrow}
\nc{\Lrarr}{\Leftrightarrow}
\rnc{\iff}{\ensuremath{\Lrarr}}
\let\le\leqslant\let\ge\geqslant
\nc{\eqn}[1]{\begin{align*}#1\end{align*}}
\nc{\dint}{\ensuremath{\displaystyle\int}}
\nc{\pd}{\partial}
\nc{\pdx}[1]{\ensuremath{\frac{\pd #1}{\pd x}}}
\nc{\pdy}[1]{\ensuremath{\frac{\pd #1}{\pd y}}}
\nc{\pdt}[1]{\ensuremath{\frac{\pd #1}{\pd t}}}
\nc{\mtc}[2]{\textcolor{#1}{\ensuremath{#2}}}
\nc{\cbox}[2]{\colorbox{#1}{\ensuremath{#2}}}
\let\tc\textcolor
\nc{\np}[1]{\ensuremath{\numprint{#1}}}
\nc{\club}{\text{\ding{168}}}
\nc{\spade}{\text{\ding{171}}}
\rnc{\diamond}{\textcolor{red}{\text{\ding{169}}}}
\nc{\heart}{\textcolor{red}{\text{\ding{170}}}}
\rnc{\bar}[1]{\overline{#1}}
\rnc{\parallel}{\mathbin{\!/\mkern-4mu/\!}}

% ──────────────────────────────────────────────
%  4 styles d'exercice — décommenter celui choisi
% ──────────────────────────────────────────────

% ---- STYLE A : texte brut, gras, proche de ton original ----
% \newcounter{exocpt}
% \newenvironment{exo}[1]{%
% 	\refstepcounter{exocpt}%
% 	\medskip
% 	\noindent\textbf{Ex.\theexocpt\ifstrempty{#1}{}{~: #1}}\par\smallskip
% }{\par}

% ---- STYLE B : numéro en couleur + filet bas ----
% \newcounter{exocpt}
% \newenvironment{exo}[1]{%
% 	\refstepcounter{exocpt}%
% 	\medskip
% 	\noindent{\bfseries\color{coulAccent}Ex.\theexocpt}%
% 	\ifstrempty{#1}{}{\enspace\textbf{:}\enspace\textbf{#1}}\\[-4pt]
% 	{\color{coulBordure}\rule{\linewidth}{0.4pt}}\par\smallskip
% }{\par}

% ---- STYLE C : ligne latérale (tcolorbox) ----
% \newcounter{exocpt}
% \newenvironment{exo}[1]{%
% 	\refstepcounter{exocpt}%
% 	\begin{tcolorbox}[
% 		enhanced, breakable, frame hidden, opacityback=0,
% 		borderline west={2pt}{0pt}{coulAccent},
% 		left=6pt, right=0pt, top=1pt, bottom=1pt,
% 		before skip=6pt, after skip=2pt,
% 		title={\textbf{Ex.\theexocpt\ifstrempty{#1}{}{~: #1}}},
% 		attach title to upper={\ ---\ },
% 		colbacktitle=white, coltitle=coulPrimaire,
% 		fonttitle=\bfseries\small,
% 	]
% }{\end{tcolorbox}}

% ---- STYLE D : cadre discret fond clair ----
\newcounter{exocpt}
\newenvironment{exo}[1]{%
	\refstepcounter{exocpt}%
	\begin{tcolorbox}[
		enhanced, breakable, arc=1pt,
		boxrule=0.4pt, colframe=coulBordure,
		colback=white,
		left=6pt, right=6pt, top=3pt, bottom=3pt,
		before skip=6pt, after skip=2pt,
		title={\textbf{Ex.\theexocpt\ifstrempty{#1}{}{~: #1}}},
		colbacktitle=coulFond, coltitle=coulPrimaire,
		fonttitle=\bfseries\small, titlerule=0pt,
		attach boxed title to top left={yshift=-1mm, xshift=3pt},
		boxed title style={boxrule=0pt, colback=coulFond},
	]
}{\end{tcolorbox}}

% ──────────────────────────────────────────────
%  Titre de série
% ──────────────────────────────────────────────
\nc{\serie}[1]{%
	\vspace{3mm}
	\noindent{\bfseries\color{coulPrimaire}\faAngleRight~#1}\\[-8pt]
	{\color{coulAccent}\rule{\linewidth}{0.6pt}}%
	\par\nopagebreak%\smallskip
}

% ──────────────────────────────────────────────
%  En-tête fiche
% ──────────────────────────────────────────────
\nc{\Entete}{
\begin{center}
    \vspace*{-7mm}
    {\LARGE\bfseries\color{coulPrimaire}\decofourleft\enspace Exercices~:~\CoursTitre\enspace\decofourright}\\[-2pt]
    {\color{coulPrimaire}\rule{\linewidth}{0.8pt}}
\end{center}
}

% ──────────────────────────────────────────────
%  Variables
% ──────────────────────────────────────────────
\nc{\CoursClasse}{2nde}
\nc{\CoursTitre}{Droites du plan}

% ══════════════════════════════════════════════
\begin{document}
\Entete

\begin{multicols}{2}

	\serie{Vecteurs directeurs et pleins de truc pas facile à traiter avec cette nouvelle fiche du exercices}

	\begin{exo}{Lecture graphique}
		Dans le repère ci-dessous, donner un vecteur directeur de chacune
		des droites $d_1$, $d_2$ et $d_3$.

		\begin{center}
			\begin{tikzpicture}[scale=0.55]
				\draw[very thin, coulBordure!40] (-0.9,-0.9) grid (5.9,5.9);
				\draw[-{Stealth}, coulPrimaire] (-0.5,0) -- (6,0)
				node[right, font=\tiny] {$x$};
				\draw[-{Stealth}, coulPrimaire] (0,-0.5) -- (0,6)
				node[above, font=\tiny] {$y$};
				\foreach \x in {1,...,5}
				\node[below, font=\tiny, coulBordure] at (\x,0) {\x};
				\foreach \y in {1,...,5}
				\node[left, font=\tiny, coulBordure] at (0,\y) {\y};
				\draw[coulAccent, ultra thick] (-0.5,-0.5) -- (5.5,5.5)
				node[above left, font=\small] {$d_1$};
				\draw[coulOrange, ultra thick] (-0.5,2) -- (5.5,2)
				node[above right, font=\small] {$d_2$};
				\draw[coulAlerte, ultra thick] (-0.5,4.5) -- (3,0)
				node[below right, font=\small] {$d_3$};
			\end{tikzpicture}
		\end{center}
	\end{exo}

	\begin{exo}{Depuis l'équation cartésienne}
		Donner un vecteur directeur de chaque droite :
		\be
		\item $d_1 : 3x - 2y + 1 = 0$
		\item $d_2 : x + 5y - 3 = 0$
		\item $d_3 : -4x + y = 0$
		\ee
	\end{exo}

	\begin{exo}{Équation depuis un point et un vecteur}
		Déterminer une équation cartésienne de $d$ passant par
		$A\coordl{2}{-3}$ et de vecteur directeur $\vec{u}\coord{4}{1}$.
	\end{exo}

	\begin{exo}{Depuis deux points}
		Déterminer une équation cartésienne de $(AB)$ avec
		$A\coordl{1}{3}$ et $B\coordl{-2}{0}$.
	\end{exo}

	\serie{Équations réduites}

	\begin{exo}{Pente et ordonnée à l'origine}
		Donner la pente $m$ et l'ordonnée à l'origine $p$ :
		\be
		\item $y = 3x - 5$
		\item $y = -\dfrac{1}{2}x + 4$
		\item $2x + 4y - 8 = 0$
		\ee
	\end{exo}

	\begin{exo}{Équation réduite depuis deux points}
		Déterminer l'équation réduite de la droite passant par
		$A\coordl{1}{5}$ et $B\coordl{3}{-1}$.
	\end{exo}
	\vfill~
	\columnbreak
	\begin{exo}{Passer d'une forme à l'autre}
		\be
		\item Donner l'équation réduite de : $$d : 6x + 3y - 5 = 0$$
		\item Donner une équation cartésienne de : $$d' : y = 6x - 5$$
		\ee
	\end{exo}

	\serie{Position relative}

	\begin{exo}{Parallèles ou sécantes ?}
		Déterminer la position relative :
		\be
		\item $d_1 : y = 2x + 1$ et $d_2 : y = 2x - 5$
		\item $d_3 : y = -x + 3$ et $d_4 : y = 4x + 1$
		\item $d_5 : 3x - y + 2 = 0$ et $d_6 : -6x + 2y - 1 = 0$
		\ee
	\end{exo}

	\begin{exo}{Démontrer le parallélisme}
		Démontrer que $d_1 : 4x - 6y + 3 = 0$ et $d_2 : -2x + 3y - 1 = 0$
		sont parallèles en utilisant les vecteurs directeurs.
	\end{exo}

	\begin{exo}{Point sur une droite}
		Les points $A\coordl{2}{7}$ et $B\coordl{3}{11}$ appartiennent-ils
		à la droite $d : y = 4x - 1$ ?
	\end{exo}

	\serie{Tracer des droites}

	\begin{exo}{Depuis l'équation cartésienne}
		Tracer la droite $d : 2x - 3y + 6 = 0$ en déterminant un point
		et un vecteur directeur.
	\end{exo}

	\begin{exo}{Plusieurs droites}
		Dans un même repère, tracer :
		\bi
		\item $d_1 : y = 2x + 1$
		\item $d_2 : y = -x + 4$
		\item $d_3 : y = 3$
		\item $d_4 : x = 2$
		\ei
	\end{exo}

	\begin{exo}{Retrouver l'équation}
		Déterminer l'équation réduite de chaque droite représentée
		ci-dessous.

		\begin{center}
			\begin{tikzpicture}[scale=0.60]
				\draw[very thin, coulBordure!40] (-0.9,-1.9) grid (5.9,5.9);
				\draw[-{Stealth}, coulPrimaire] (-0.5,0) -- (6,0)
				node[right, font=\tiny] {$x$};
				\draw[-{Stealth}, coulPrimaire] (0,-1.5) -- (0,6)
				node[above, font=\tiny] {$y$};
				\foreach \x in {1,...,5}
				\node[below, font=\tiny, coulBordure] at (\x,0) {\x};
				\foreach \y in {-1,...,5}
				\node[left, font=\tiny, coulBordure] at (0,\y) {\y};
				\draw[coulAccent, ultra thick, domain=-0.5:4.5]
				plot(\x, {\x+1}) node[right, font=\small] {$d_1$};
				\draw[coulOrange, ultra thick, domain=-0.5:2.5]
				plot(\x, {-2*\x+4}) node[right, font=\small] {$d_2$};
				\filldraw[coulAccent] (0,1) circle (0.1);
				\filldraw[coulAccent] (2,3) circle (0.1);
				\filldraw[coulOrange] (0,4) circle (0.1);
				\filldraw[coulOrange] (2,0) circle (0.1);
			\end{tikzpicture}
		\end{center}
	\end{exo}
	% \vfill~\columnbreak
	\columnbreak
	\begin{exo}{Retrouver l'équation}
		Déterminer l'équation réduite de chaque droite représentée
		ci-dessous.

		\begin{center}
			\begin{tikzpicture}[scale=0.65]
				\draw[very thin, coulBordure!40] (-0.9,-1.9) grid (5.9,5.9);
				\draw[-{Stealth}, coulPrimaire] (-0.5,0) -- (6,0)
				node[right, font=\tiny] {$x$};
				\draw[-{Stealth}, coulPrimaire] (0,-1.5) -- (0,6)
				node[above, font=\tiny] {$y$};
				\foreach \x in {1,...,5}
				\node[below, font=\tiny, coulBordure] at (\x,0) {\x};
				\foreach \y in {-1,...,5}
				\node[left, font=\tiny, coulBordure] at (0,\y) {\y};
				\draw[coulAccent, ultra thick, domain=-0.5:4.5]
				plot(\x, {\x+1}) node[right, font=\small] {$d_1$};
				\draw[coulOrange, ultra thick, domain=-0.5:2.5]
				plot(\x, {-2*\x+4}) node[right, font=\small] {$d_2$};
				\filldraw[coulAccent] (0,1) circle (0.1);
				\filldraw[coulAccent] (2,3) circle (0.1);
				\filldraw[coulOrange] (0,4) circle (0.1);
				\filldraw[coulOrange] (2,0) circle (0.1);
			\end{tikzpicture}
		\end{center}
	\end{exo}
	\serie{Tracer des droites}

	\begin{exo}{Depuis l'équation cartésienne}
		Tracer la droite $d : 2x - 3y + 6 = 0$ en déterminant un point
		et un vecteur directeur.
	\end{exo}

	\begin{exo}{Plusieurs droites}
		Dans un même repère, tracer :
		\bi
		\item $d_1 : y = 2x + 1$
		\item $d_2 : y = -x + 4$
		\item $d_3 : y = 3$
		\item $d_4 : x = 2$
		\item $d_5 : x = 2 - y$
		\ei
	\end{exo}
	\vfill~\columnbreak
	\begin{exo}{}
		Prout

	\end{exo}
	\begin{exo}{}
		Les points $I$, $J$ et $K$ sont les milieux des côtés $[AB]$, $[AC]$ et $[BC]$.

		\begin{center}
			\begin{tikzpicture}[scale=0.7]
				% Grille
				\draw[very thin, coulBordure!40] (-4.9,-3.9) grid (5.9,2.9);
				% Axes
				\draw[-{Stealth}, coulPrimaire] (-5,0) -- (6,0);
				\draw[-{Stealth}, coulPrimaire] (0,-4) -- (0,3);
				% Vecteurs i et j
				\draw[-{Stealth[length=5pt]}, coulPrimaire, thick] (0,0) -- (1,0);
				% node[below, font=\small] {$\vec{i}$};
				\draw[-{Stealth[length=5pt]}, coulPrimaire, thick] (0,0) -- (0,1);
				% node[right, font=\small] {$\vec{j}$};
				% Points
				\coordinate (A) at (-4,2);
				\coordinate (B) at (4,-2);
				\coordinate (C) at (2,2);
				\coordinate (I) at (0,0);   % milieu AB
				\coordinate (J) at (-1,2);   % milieu AC
				\coordinate (K) at (3,0); % milieu BC — ajusté à la grille
				% Triangle rempli
				\fill[tealblue, opacity=0.15] (A) -- (B) -- (C) -- cycle;
				\draw[tealblue, thick] (A) -- (B) -- (C) -- cycle;
				% Points nommés
				\foreach \P/\nom/\pos in {
						I/I/below left,
						J/J/above,
						K/K/above right}{
						\filldraw[coulOrange] (\P) circle (0.08)
						node[\pos, font=\small, coulOrange] {$\nom$};
					}
				% Repasser A, B, C en teal
				\foreach \P/\nom/\pos in {A/A/above left, B/B/below right, C/C/above right}{
						\filldraw[tealblue] (\P) circle (0.08)
						node[\pos, font=\small, tealblue] {$\nom$};
					}
			\end{tikzpicture}
		\end{center}

		\textbf{1.} Déterminer les équations réduites des droites $(AK)$ et $(BJ)$.

		\textbf{2.} En déduire les coordonnées de leur point d'intersection $G$.

		\textbf{3.} Démontrer que $G$ appartient également à la droite $(CI)$.
	\end{exo}
\end{multicols}

\end{document}
